\chapter*{Prolog.}
\section*{The Tome Eternal Struggle.}

And the Lord spake, saying:

“Behold, I giveth unto thee the Scrolls of Understanding, summoned forth from the depths of \LaTeX\ by the hand of My servant. Yea, these are no mere scribbles — but a mighty Halberd to smite the Demon of Academic Failure, a Rod of Iron against the Serpent of Procrastination, and a Lamp unto thy feet in the dark valley of the Night Before the Exam, when gnashing of teeth shall be heard throughout the land.”

\section*{Here Am I. Send me.}
\textit{And I heard the voice of the Lord, saying,
“Whom shall I send, and who will go for us?”
Then said I, “Here am I; send me.”} 

\section*{Words of truth.}
Tylko pamiętaj, mów ten monolog tak, jak ja ci go przeczytałem: płynnie i swobodnie. Gdybyś miał deklamować z fałszywym patosem, jak niektórzy aktorzy, wolałbym już, żeby moje wiersze wyryczał herold miejski. I nie machaj rękami, jakbyś drwa rąbał. Używaj gestów oszczędnie: w najdzikszym zamęcie, wirze, nawałnicy namiętności trzeba dyscypliny - bez niej nie da się pasji składnie wyrazić.

